% GENERATED FILE. Edit canonical lessons, then run npm run generate:books.
% canonical-source-hash: fnv1a64:28d14a23010faa47
% canonical-lessons: PA-R48-forms-name-and-language, PA-R48-forms-residence-and-work, PA-R48-forms-cold-all-nine, PA-R48-gurmukhi-met-once, PA-R48-marks-above-and-below, PA-R48-courtesy-and-parting

\chapter{Second Pass --- What Was Met Once}
\label{ch:pa-doosra-gera}
\hlchaptermodality{48}

\begin{chapteropening}
\textbf{By the end of this chapter:} \emph{I can fill every A1 form field at each level of support, read every Gurmukhi piece this book taught once, and place each courtesy and parting word in its register.}

From cold: say please, sorry and welcome, name the verb or case ending inside each, and place the three partings in their registers.
\end{chapteropening}

\section[nāṁ · bhāshā]{\pa{ਨਾਂ} and \pa{ਭਾਸ਼ਾ} — the same two fields, three times over}

\label{lesson:PA-R48-forms-name-and-language}



\begin{quote}
The practice name this book gave you to write. Say it, then say what label stands to the left of it on a form.
\end{quote}

\begin{culture}
Each field in this book was taught three times, and the three are not repetition. They are three different amounts of help:

\noindent
\begin{tabularx}{\linewidth}{@{}>{\raggedright\arraybackslash}X>{\raggedright\arraybackslash}X>{\raggedright\arraybackslash}X@{}}
\toprule
\textbf{stage} & \textbf{what is on the page} & \textbf{what it proves} \\
\midrule
\textbf{supported} & a visible bank of values to copy from & you can place a value in a field \\
\textbf{delayed} & the bank is shown, then covered & you can hold it while you write \\
\textbf{no model} & a label, a blank, and nothing else & you can produce it \\
\bottomrule
\end{tabularx}

Only the third is evidence that you can write the thing. The first two are evidence that you can \emph{copy} and \emph{remember} — which are real skills, and different ones. A form you can only fill with the answer visible is a form you cannot fill.
\end{culture}

\subsection*{Your turn}
Cold, with nothing in front of you.

\begin{itemize}
\raggedright
  \item \emph{Write it:} the name label, then a name you know
  \item \emph{Write it:} the language label, then the language this book is in
  \item \emph{Say it:} the three field labels the book asks you to join in one pass
\end{itemize}

\begin{tcolorbox}[breakable,colback=teal!4,colframe=teal!35!black,title={Before you move on}]
Which of the three stages is the only one that proves you can write a field? (\textbf{The one with no model.}) What is on the page at that stage? (\textbf{A label and a blank.})
\end{tcolorbox}
\section[rihāish · kamm]{\pa{ਰਿਹਾਇਸ਼} and \pa{ਕੰਮ} — and the move for when a field is already wrong}

\label{lesson:PA-R48-forms-residence-and-work}



\begin{quote}
Two more labels. Say the one that asks where you live, then the one that asks what you do.
\end{quote}

\subsection*{Your turn: the same ladder, two more fields}
Write these out:

\begin{itemize}
\raggedright
  \item the residence label and a value, with nothing to copy
  \item the work label and a value, with nothing to copy
\end{itemize}

Both fields climbed the same three rungs the name and language fields did. What changes between fields is the vocabulary, not the ladder.

\begin{culture}
\textbf{Repair is its own skill, and only the residence field teaches it.} The move is not \emph{fill the field} — it is \emph{one dimension of a completed field is wrong, change that dimension and leave the rest standing}.

That is a harder thing than it sounds, and it is the thing a real form actually demands: a clerk points at one line, you correct one line. A learner who can only write a field from scratch will rewrite the whole row and introduce two new errors doing it.
\end{culture}

\begin{tcolorbox}[breakable,colback=teal!4,colframe=teal!35!black,title={Before you move on}]
When you repair a field, how much of it do you rewrite? (\textbf{Only the dimension that is wrong.}) Which field in this book taught repair? (\textbf{Residence.})
\end{tcolorbox}
\section[nāṁ · bhāshā · rihāish]{Nine moves, sorted by how much the page is helping}

\label{lesson:PA-R48-forms-cold-all-nine}



\begin{quote}
Two fields, no model, nothing on the page but a label and a blank. Write both before reading on.
\end{quote}

\subsection*{Your turn: across the three fields, one rung at a time}
The book taught these field by field. Take them rung by rung instead.

Write these out:

\begin{itemize}
\raggedright
  \item all three fields with a bank in front of you — the \textbf{supported} rung
  \item all three after the bank is covered — the \textbf{delayed} rung
  \item all three from a label alone — the \textbf{no-model} rung
\end{itemize}

\begin{grammarlens}[title={What you've built}]
Read across rather than down and the ladder stops being about forms at all. It is about \textbf{what is in the room when you produce something}: a page that shows you the answer, a page that showed you and stopped, and a page that never did.

Three fields multiplied by three rungs is nine separate things you can now do, and the only three that count as writing evidence are the bottom row. The other six are the scaffolding that got you there, and scaffolding is worth naming rather than mistaking for the building.
\end{grammarlens}

\begin{tcolorbox}[breakable,colback=teal!4,colframe=teal!35!black,title={Before you move on}]
Of the nine moves, how many are writing evidence? (\textbf{Three} — the no-model row.) What is the difference between the supported and delayed rungs? (\textbf{The bank is visible, then covered.})
\end{tcolorbox}
\section[īṛī · iṛī · ūṛā · ṭhaṭṭhā]{\pa{ਈ} \pa{ਇ} \pa{ਉ} \pa{ਠ} — the pieces that were met once}

\label{lesson:PA-R48-gurmukhi-met-once}



\begin{quote}
Two standing vowels. Read \textbf{\pa{ਇ}}, then \textbf{\pa{ਉ}}, and say which vowel each one carries when nothing precedes it.
\end{quote}

\subsection*{Script: four standing shapes and one mark}
\noindent
\begin{tabularx}{\linewidth}{@{}>{\raggedright\arraybackslash}X>{\raggedright\arraybackslash}X>{\raggedright\arraybackslash}X@{}}
\toprule
\textbf{shape} & \textbf{name} & \textbf{what it is} \\
\midrule
\textbf{\pa{ਈ}} & īṛī & a standing long \textbf{ī} \\
\textbf{\pa{ਇ}} & iṛī & a standing short \textbf{i} \\
\textbf{\pa{ਉ}} & ūṛā & a standing \textbf{u} \\
\textbf{\pa{ਠ}} & ṭhaṭṭhā & a consonant — the curled retroflex \\
\textbf{\pa{ੋ}} & horā & a \textbf{mark}, not a letter: long \textbf{o}, written above \\
\bottomrule
\end{tabularx}

The first three are the ones that sit at the front of a word with no consonant to lean on. The last is the opposite kind of thing entirely — it cannot stand anywhere, and it goes above the letter it belongs to.

\subsection*{Your turn: read them where you met them}
Whole words this time, each read cold.

Read these aloud:

\begin{itemize}
\raggedright
  \item \textbf{\pa{ਨਮਸਤੇ}} — the greeting, five pieces, all of them yours
  \item \textbf{\pa{ਚੰਗਾ}} — the everyday sign-off word
  \item \textbf{\pa{ਠੀਕ}-\pa{ਠਾਕ}} — and find the \textbf{\pa{ਠ}} twice
  \item \textbf{\pa{ਜੀ} \pa{ਆਇਆਂ} \pa{ਨੂੰ}} — the welcome
  \item \textbf{\pa{ਸ਼ੁਭ} \pa{ਰਾਤ}} — the formal good night
\end{itemize}

\begin{tcolorbox}[breakable,colback=teal!4,colframe=teal!35!black,title={Before you move on}]
Which of the five shapes above cannot stand on its own? (\textbf{\pa{ੋ}}, the horā — it is a mark.) Which word in the practice carries \textbf{\pa{ਠ}} twice? (\textbf{\pa{ਠੀਕ}-\pa{ਠਾਕ}}.)
\end{tcolorbox}
\section[āṛā · dulāvāṁ · auṅkaṛ]{\pa{ਆ} \pa{ੈ} \pa{ੁ} — one that stands, two that ride}

\label{lesson:PA-R48-marks-above-and-below}



\begin{quote}
The standing long \textbf{ā}. Read it, and say which plain letter it is built on before the mark was added.
\end{quote}

\subsection*{Script: where a mark sits is what it means}
\noindent
\begin{tabularx}{\linewidth}{@{}>{\raggedright\arraybackslash}X>{\raggedright\arraybackslash}X>{\raggedright\arraybackslash}X@{}}
\toprule
\textbf{shape} & \textbf{name} & \textbf{where it sits} \\
\midrule
\textbf{\pa{ਆ}} & āṛā & nowhere — it \textbf{is} the letter \\
\textbf{\pa{ੈ}} & dulāvāṁ & \textbf{above}, two strokes: \emph{ai} \\
\textbf{\pa{ੁ}} & auṅkaṛ & \textbf{below}, a hook: short \emph{u} \\
\bottomrule
\end{tabularx}

Gurmukhi puts vowel marks above, below, before and after the consonant they attach to, and the position is not decoration — it is how you tell one mark from another at a glance. Two strokes above is \emph{ai}; a hook below is \emph{u}. Change the position and you have changed the vowel.

\textbf{\pa{ਆ}} is the counter-example that makes the rule visible: it is a full letter wearing a mark that has fused into it, so it stands where a consonant would rather than hanging off one.

\subsection*{Your turn}
\begin{itemize}
\raggedright
  \item \emph{Read it:} \textbf{\pa{ਮੈਨੂੰ}} — and find the two strokes above
  \item \emph{Read it:} \textbf{\pa{ਜੀ} \pa{ਆਇਆਂ} \pa{ਨੂੰ}} — the welcome, opening on the standing \textbf{\pa{ਆ}}
  \item \emph{Say it:} the formal good night, and say which parting it stands beside
\end{itemize}

\textbf{\pa{ਸ਼ੁਭ} \pa{ਰਾਤ}} is the Sanskritic one. It is not the everyday sign-off; reaching for it among friends sounds like reading from a card.

\begin{tcolorbox}[breakable,colback=teal!4,colframe=teal!35!black,title={Before you move on}]
Which mark sits above, and which below? (\textbf{Dulāvāṁ} above, \emph{ai}; \textbf{auṅkaṛ} below, short \emph{u}.) Is \textbf{\pa{ਸ਼ੁਭ} \pa{ਰਾਤ}} the formal parting or the everyday one? (\textbf{Formal.})
\end{tcolorbox}
\section[miharbānī karke · māf karo · rabb rākhā]{\pa{ਮਿਹਰਬਾਨੀ} \pa{ਕਰਕੇ} — the politeness words, and what is inside them}

\label{lesson:PA-R48-courtesy-and-parting}



\begin{quote}
Say \emph{please}. Then say \emph{sorry}. Both are longer than the English and both have a verb inside — say which verb.
\end{quote}

\begin{culture}
Punjabi's two commonest courtesies are not words for \emph{please} and \emph{sorry} at all. They are \textbf{things done}:

\noindent
\begin{tabularx}{\linewidth}{@{}>{\raggedright\arraybackslash}X>{\raggedright\arraybackslash}X@{}}
\toprule
\textbf{what you say} & \textbf{what it is built from} \\
\midrule
\textbf{\pa{ਮਿਹਰਬਾਨੀ} \pa{ਕਰਕੇ}} & \emph{having done a kindness} \\
\textbf{\pa{ਮਾਫ਼} \pa{ਕਰੋ}} & \emph{make pardoned} \\
\textbf{\pa{ਜੀ} \pa{ਆਇਆਂ} \pa{ਨੂੰ}} & \emph{to those who have come — life} \\
\bottomrule
\end{tabularx}

All three run on the same verb of doing or a case ending, and none of them is a fixed politeness particle you could swap for another. The \textbf{\pa{ਨੂੰ}} on the third is the same dative ending the book uses elsewhere: the welcome is grammatically \emph{addressed} to the arrivals rather than describing them.
\end{culture}

\subsection*{Your turn: leaving, and the registers}
Say these aloud:

\begin{itemize}
\raggedright
  \item we will meet — and say what the \textbf{-\pa{ਾਂਗੇ}} ending is doing
  \item the parting that hands someone to God's keeping — \emph{rabb rākhā}
  \item the formal, Sanskritic good night — \emph{shubh rāt}
  \item the bare negative — \emph{nā}
\end{itemize}

Three partings, three registers. \emph{Rabb rākhā} carries a blessing and belongs to people who know each other; \emph{shubh rāt} is the formal one; and the everyday sign-off is neither.

\begin{tcolorbox}[breakable,colback=teal!4,colframe=teal!35!black,title={Before you move on}]
What is literally inside the word for \emph{please}? (\textbf{Having done a kindness.}) What does the echo in \textbf{\pa{ਠੀਕ}-\pa{ਠਾਕ}} add — a new meaning, or a shade of one? (\textbf{A shade} — the echo softens it toward \emph{getting by}.) And which body part is a cousin of English \emph{horn} rather than \emph{head}? (\textbf{\pa{ਸਿਰ}}.)
\end{tcolorbox}
